\chapter{Các công trình nghiên cứu}
\label{ch:congtrinh}

Nội dung của đề tài được viết thành ba bài báo tách rời, khác nhau về đóng góp chính, về
ngôn ngữ và về số liệu sử dụng.

\begin{quote}
Lê Đoàn Phương Uyên, \GVHD. ``Descriptor and Preference Targets for GUI Instruction
Generation''.
\end{quote}

\noindent
Bản thảo tiếng Anh, độ dài $7$ trang theo mẫu IEEE, gửi ngày $31$ tháng $8$ năm $2026$ tới
Hội thảo khoa học quốc gia \emph{Nghiên cứu cơ bản và ứng dụng Công nghệ thông tin} (FAIR
2026), tiểu ban Xử lý ngôn ngữ tự nhiên. Đóng góp chính là phần \textbf{mô hình}, gồm thiết kế so sánh ở
Chương~\ref{ch:phuongphap}, kết quả của các nhánh ở Chương~\ref{ch:thucnghiem}, phép tách
công giữa mục tiêu ưu tiên và nhánh so sánh, và chẩn đoán điểm nghẽn ở
Mục~\ref{sec:chonghen}. Thước đo xuất hiện trong bài với vai trò \emph{dụng cụ đã hiệu chuẩn},
được trình bày gọn trong một mục, chứ không phải với vai trò một đóng góp.

\begin{quote}
Lê Đoàn Phương Uyên, \GVHD. ``Sinh hướng dẫn sử dụng phần mềm từ ảnh chụp màn hình và mục
tiêu người dùng: xây dựng nhãn quy chiếu tự động khi phần tử giao diện không có tên để
gọi''.
\end{quote}

\noindent
Bản thảo tiếng Việt, nộp ngày $30$ tháng $8$ năm $2026$ tại \emph{Hội thảo Quốc gia lần thứ
tư về Ngôn ngữ học Tính toán} (VCL 2026). Đóng góp chính là phần \textbf{nhãn quy chiếu và quy trình dựng dữ liệu} ở
Chương~\ref{ch:dulieu}. Bài toán được đặt trong dòng sinh biểu thức quy chiếu rồi áp vào
miền giao diện, trong đó $22{,}0\%$ phần tử cần chạm không có tên nào để gọi và $7{,}6\%$
trùng tên với
một phần tử khác trên cùng màn hình.

\begin{quote}
Lê Đoàn Phương Uyên, \GVHD. ``Separating the Action Gate from Geometry in Execution-Based
Evaluation for GUI Instruction Generation''.
\end{quote}

\noindent
Bản thảo tiếng Anh theo mẫu Springer LNCS, nộp ngày $21$ tháng $9$ năm $2026$ tới hội nghị
SOICT 2026 (mã bài $5577$), kỷ yếu dự kiến thuộc tùng thư CCIS. Đóng góp chính là một điểm về
đo lường: khi chấm câu hướng dẫn bằng thước dựa trên thực thi, phần mà điều kiện loại thao
tác đóng góp phải được tách khỏi phần định vị, vì mức hơn của chặng ba so với nhánh chỉ sinh câu
dưới luật khớp chạm của AndroidInTheWild đến gần như trọn từ nhóm bước nhánh nền gọi sai loại
thao tác (Mục~\ref{sec:nhieuthuoc}).

Hai bài đầu chia số liệu theo một bảng phân chia cố định để không chồng lấn. Phần thước đo và
toàn bộ điểm số của các nhánh thuộc bài thứ nhất. Phần kiểm phép ghép dữ liệu, chất lượng
nhãn và hai bộ lọc thuộc bài thứ hai. Quy mô dữ liệu là phần chung, được mô tả đầy đủ ở bài
thứ hai và chỉ được nêu trong một câu ở bài thứ nhất. Vì hai bài nộp cách nhau một ngày, phần trích chéo
giữa chúng ghi ở dạng \emph{đang bình duyệt}. Nhánh ứng viên ở Mục~\ref{sec:nhanhungvien}
và các nhánh từ cuối tháng $9$ (Mục~\ref{sec:kenhvitri} đến Mục~\ref{sec:tage}) không thuộc bài
nào trong ba bài, vì chúng được thực hiện sau ngày nộp.
